% Abstracts are usually written in English, with a version in your
% mother tongue underneath
\chapter*{Abstract} 
\addcontentsline{toc}{chapter}{Abstract}
% Don't change anything above this.
In late November 2025, Sumatra was devastated by the rare equatorial Tropical Cyclone (TC) Senyar, which triggered catastrophic flash floods and landslides affecting millions of people. While rapid attribution studies suggested a human fingerprint, these unconditional approaches might not able to capture the local and nonlinear processes, unique to the event. This study utilizes a high-resolution, convection-permitting regional climate model (RegCM5) and a pseudo-global warming (PGW) storyline approach to attribute the influences of anthropogenic climate change (ACC) to the event and project future evolution under a +1.5 K warming scenario.

Reconstruction of the event confirms that the convection-permitting model significantly improves the representation of the storm's spatial core and intensity compared to reanalysis data. Historical attribution reveals that ACC has already amplified the disaster: while total rainfall in Aceh increased by 10\%, the rainfall area exceeding 300 mm expanded by a massive 78\%, drastically increasing geographic exposure to hydrological triggers. Future projections indicate a more violent, concentrated bursts response: while the total rainfall and area are projected to contract by 9\% and 14\% respectively, heavy rainfall intensity continues to increase (+8\%), exceeding Clausius-Clapeyron scaling.

These shifts are underpinned by a robust increase in atmospheric energy and moisture transport, including a 24\% increase in CAPE and an intensification of extreme moisture transport (+7–10\% per K). The increased release of latent heat within the storm core accelerates updraft velocities, effectively converting moisture to extreme precipitation. Furthermore, while TC Senyar’s trajectory remains robust across climate states, a consistent increase in peak wind intensity in future scenarios suggests an emerging threat of compounding wind and hydrological hazards for equatorial Sumatra. These findings imply that future adaptation strategies must prioritize resilience against high-magnitude, localized bursts even within a potentially smaller total disaster footprint.

% A short, abstracted description of your research project goes here. 
% It should be about 100 words long. But write it last.

% An abstract is not a summary of your research project: it's an abstraction of that. 
% It tells the readers why they should be interested in your research project but summarises all
% they need to know if they read no further.

% The writing style used in an abstract is like the style used in the rest of your research project: concise, clear and direct. 
% In the rest of the research project, however, you will introduce and use technical terms. In the abstract you should
% avoid them in order to make the result comprehensible to all.

% You may like to repeat the abstract in your mother tongue.

% At a unviersity like Stellenbosch you *must* produce an abstract in Afrikaans for your masters.
% At AIMS you are encouraged to repeat the abstract in your mother tongue
% French, Igbo, Mlagasy, etc. just write it using LaTeX's special
% characters.
% Arabic students see the arabic.tex file for an example
% Amharic use openoffice and export from there and import a figure here.
% Where the words do not exist put the English work in italics, or use mathematical symbols.


% Do not change anything below this except for adding your
% signature (replace images/signature.png) and your name.
\vfill
\section*{Declaration}
I, the undersigned, hereby declare that the work contained in this research project is my original work, and that any work done by others or by myself previously has been acknowledged and referenced accordingly.

% Scan your signature into a small picture called 'signature.png' and insert it
% above your name and the date:
\hspace{1cm}
%  Sign here
\includegraphics[height=3cm]{figs/sign.png} \newline \hrule
% Your name must be in English Capitalisation with no comma, and the Family name comes last. 
% Do note the date below. It is called the "deadline".
Nathasya Abenita Christien
\newline
August 20, 2025

